% !TeX spellcheck = en_US
\addsection{Resources}{\skills/estates.png}

\iftoggle{printable}{\vspace{-\baselineskip}}{}

\begin{multicols}{2}

\pagetarget{Resources}{There}\index{Resources} are three types of Resources in the game: Gold \svg{gold}, Building Materials \svg{building_materials}, and Valuables \svg{valuablegreater}.
Resources are spent during the game to expand your \pagelink{Town}{Town}, to Recruit \pagelink{Units}{Units}, and to purchase \pagelink{spells}{Spells}.
You can gain Resources from the \pagelink{Mines}{Settlements and Mines} that you have \pagelink{Categories}{Flagged}, and also by playing Cards and rolling Resource Dice \svg{resource_die}.
Whenever a player's Resource Production is increased or decreased, move that Resource's cube on its production track the appropriate number of spaces.\par
\begin{multicols}{3}
  \centering
  \vspace*{\fill}
  \includegraphics[width=0.8\linewidth]{\images/gold.png}
  \footnotesize\textcolor{darkcandyapplered}{\textit{\textbf{Gold \phantom{Materials}}}}
  \columnbreak
  \vspace*{\fill}
  \includegraphics[width=0.8\linewidth]{\images/building_materials.png}
  \footnotesize\textcolor{darkcandyapplered}{\textit{\textbf{Building Materials}}}
  \columnbreak
  \vspace*{\fill}
  \includegraphics[width=0.8\linewidth]{\images/valuables.png}
  \footnotesize\textcolor{darkcandyapplered}{\textit{\textbf{Valuables \phantom{Materials}}}}
\end{multicols}
Players start each Scenario with the number of Resources indicated in that Scenario’s setup.
Resources can also be \pagelink{Trading}{traded}.
There's no limit to the amount of Resources you can have.

\vspace*{\fill}

\columnbreak

\begin{center}
  \includegraphics[width=\linewidth]{\tables/resource_board.png}
  \footnotesize\textcolor{darkcandyapplered}{\textbf{\textit{Resource Production Tracker}}}
\end{center}
\bigskip

Possible Resource Die\index{Resource Die}\index{Dice} \svg{resource_die} results:
\medskip
\begin{itemize}
  \setlength\itemsep{8pt}
  \item[{\svg[15]{2_building_materials}}] — 2 × Building Materials
  \item[{\svg[15]{4_building_materials}}] — 4 × Building Materials
  \item[{\svg[15]{1_valuables}}] — 1 × Valuables
  \item[{\svg[15]{2_valuables}}] — 2 × Valuables
  \item[{\svg[15]{3_gold}}] — 3 × Gold
  \item[{\svg[15]{6_gold}}] — 6 × Gold
\end{itemize}

\begin{tikzpicture}[overlay]
  \node at (6, 1.5) {\includegraphics[width=0.3\linewidth]{\images/resource_die.png}};
\end{tikzpicture}

\end{multicols}

\vspace*{\fill}

\begin{scaledfigure}[blanker]
  \centering
  \includegraphics[width=\linewidth, height=\myspace, keepaspectratio]{\art/griffin.jpg}
\end{scaledfigure}
